\clearpage
\section{Marco Teórico}

Para construir el marco teórico es necesario que la información que se
reúna y analice, se relacione de manera coherente entre sí y con el objeto
de estudio, evitando divagar con el planteamiento de temas no vinculados
directamente al objeto de investigación. Se incluirá toda aquella
información ya documentada que se tiene sobre el tema de estudio,
realizando una sistematización del conocimiento científico en relación con
la temática a estudiar, que evidencie el avance de la ciencia y los vacíos
de conocimientos que orientan la investigación. Toda la información
bibliográfica usada debe estar respaldada por sus respectivas citas, las
cuales deben ser incluidas en la lista de referencias bibliográficas en
formato IEEE.

% ============================================================
\subsection{Fundamentos teóricos}
% ============================================================

% ------------------------------------------------------------
\subsubsection{Gestión Nutricional Clínica}
% ------------------------------------------------------------

La gestión nutricional en centros especializados ecuatorianos sigue siendo
predominantemente manual, apoyada en hojas de cálculo o registros físicos.
Este modelo genera vulnerabilidad en los datos, ralentiza la atención y
dificulta el acceso rápido al historial del paciente~\cite{quichimbo2021}.
La falta de digitalización incrementa los errores en la planificación
dietética y debilita el vínculo entre el profesional y el paciente, lo que
evidencia la necesidad de herramientas tecnológicas que complementen la
consulta clínica sin desplazar el juicio del
nutricionista~\cite{vargas2025}.

% ------------------------------------------------------------
\subsubsection{Nutrición Personalizada}
% ------------------------------------------------------------

La Nutrición Personalizada~(NP) propone utilizar información específica
del paciente (patológica, fisiológica y hereditaria) para diseñar
intervenciones dietéticas más efectivas que las guías
generales~\cite{siddiqui2025}. Estudios recientes evidencian que la
integración de catálogos alimentarios adaptados al contexto local resulta
esencial para garantizar la pertinencia clínica de las
recomendaciones~\cite{aldarras2025}. A diferencia de los enfoques
genéricos, la NP considera variantes metabólicas individuales, lo que
exige plataformas capaces de procesar datos heterogéneos del paciente de
forma eficiente~\cite{singar2024}.

\begin{figure}[ht]
  \centering
  \includegraphics[width=0.80\linewidth]{imagenes/Imagen3.jpg}
  \caption{Imagen de ``Sistema de Nutrición Personalizada (NP)'' generada por Gemini.}
  \label{fig:np}
\end{figure}

% ------------------------------------------------------------
\subsubsection{Sistemas de Soporte a la Decisión Clínica}
% ------------------------------------------------------------

Los Sistemas de Soporte a la Decisión Clínica~(SSDC) permiten integrar y
procesar datos nutricionales de forma ágil, facilitando la toma de
decisiones del profesional de salud y reduciendo errores en la
planificación dietética~\cite{quichimbo2021}. Las plataformas de salud
móvil~(mHealth) fortalecen la interacción entre paciente y especialista,
promoviendo un seguimiento continuo y mayor adherencia al plan
nutricional~\cite{ghose2021}. Sin embargo, estudios sistemáticos
evidencian errores de hasta un~22\% en la predicción de macronutrientes en
plataformas actuales, lo que limita su aplicación clínica sin supervisión
profesional~\cite{mortazavi2023}, evidenciando vacíos que el presente
proyecto busca abordar.

% ------------------------------------------------------------
\subsubsection{Modelos de Lenguaje de Gran Escala y RAG}
% ------------------------------------------------------------

Los modelos de lenguaje de gran escala~(LLM) son sistemas de inteligencia
artificial basados en arquitecturas de transformadores que aprenden
representaciones lingüísticas complejas, permitiéndoles generar respuestas
coherentes ante consultas en lenguaje natural. No obstante, los LLM pueden
producir alucinaciones al no disponer de un contexto verificado. La
Generación Aumentada por Recuperación~(RAG) resuelve esta limitación
combinando un módulo de recuperación de información con el LLM: ante una
consulta, el sistema recupera los fragmentos más relevantes de una base de
conocimiento estructurada y los entrega al modelo como contexto,
reduciendo las alucinaciones al anclar las respuestas en datos verificados
y específicos del dominio~\cite{gavai2025,arshad2025}. En sistemas de
nutrición personalizada, esta arquitectura ha reportado niveles de
adherencia de hasta un~80\% y la capacidad de ajustar dinámicamente los
planes generados~\cite{gavai2025}.

% ============================================================
\subsection{Modelos, métodos y técnicas}
% ============================================================

% ------------------------------------------------------------
\subsubsection{Pipeline RAG y Almacenamiento Vectorial}
% ------------------------------------------------------------

El pipeline RAG del prototipo sigue tres etapas secuenciales. En la
indexación, los documentos del catálogo alimentario ecuatoriano
y los perfiles clínicos básicos del paciente se convierten en vectores
numéricos~(embeddings) mediante un modelo de embeddings de código abierto
y se almacenan en pgvector, extensión de~PostgreSQL que permite
búsquedas de similitud semántica sobre datos relacionales. En la etapa de
recuperación, ante una consulta del nutricionista, el sistema
calcula la similitud vectorial y retorna los fragmentos más relevantes.
En la generación, el LLM produce la recomendación nutricional
usando esos fragmentos como contexto. LangChain actúa como framework de
orquestación del pipeline completo.

\begin{figure}[ht]
  \centering
  \includegraphics[width=0.95\linewidth]{imagenes/Imagen1.jpg}
  \caption{Imagen de ``Pipeline RAG del prototipo NutriRAG'' generada por Gemini.}
  \label{fig:pipeline_rag}
\end{figure}

% ------------------------------------------------------------
\subsubsection{Modelos LLM de Código Abierto}
% ------------------------------------------------------------

Para el pipeline se evalúan modelos de lenguaje de código abierto con
soporte multilingüe ejecutables sobre infraestructura local. Entre los
candidatos se consideran: \textbf{Qwen}~(Alibaba Cloud), con sólido
soporte multilingüe y seguimiento de instrucciones;
\textbf{Kimi~K}~(Moonshot~AI), destacado por su manejo de contextos
largos; \textbf{GLM}~(Zhipu~AI), con capacidad de razonamiento en
entornos técnicos; y \textbf{Nemotron}~(NVIDIA), optimizado para
seguimiento preciso de instrucciones en dominios
especializados~\cite{gavai2025,arshad2025}. El uso de modelos locales
garantiza que los datos clínicos del paciente no sean transmitidos a
servicios externos, respondiendo a las limitaciones de privacidad
identificadas en la literatura~\cite{abeltino2025,ebadi2025}.

% ------------------------------------------------------------
\subsubsection{Escala de Satisfacción tipo Likert y Análisis Descriptivo}
% ------------------------------------------------------------

La escala Likert es un instrumento ampliamente
utilizado para medir actitudes y percepciones en investigación en
sistemas de información~\cite{ghose2021}. En este proyecto se emplea una
escala de cinco puntos (1~=~Muy insatisfecho; 5~=~Muy satisfecho) que
evalúa cinco dimensiones: adecuación clínica, claridad del lenguaje,
facilidad de uso, confianza en la información generada y eficiencia
temporal percibida. El índice global de satisfacción se calcula como el
promedio ponderado de los ítems por dimensión, y los resultados se
analizan mediante estadística descriptiva~(media, desviación estándar y
distribución de frecuencias).

\begin{figure}[ht]
  \centering
  \includegraphics[width=0.85\linewidth]{imagenes/Imagen4.jpg}
  \caption{Imagen de ``Instrumento de satisfacción tipo Likert del sistema NutriRAG'' generada por Gemini.}
  \label{fig:likert}
\end{figure}

% ------------------------------------------------------------
\subsubsection{Stack Tecnológico y Modelo C4}
% ------------------------------------------------------------

El prototipo se desarrolla con un stack de código abierto: backend en
Node.js con TypeScript y Express.js, frontend en React con TailwindCSS,
y persistencia en PostgreSQL con Prisma~ORM y la extensión pgvector. La
arquitectura del sistema se documenta mediante el modelo~C4~(Contexto,
Contenedores, Componentes y Código), que proporciona una descripción
estructurada de los módulos, sus responsabilidades y sus interacciones,
constituyendo la guía de implementación del prototipo.

\begin{figure}[ht]
  \centering
  \includegraphics[width=0.75\linewidth]{imagenes/Imagen2.jpg}
  \caption{Imagen de ``Diagrama C4 de contexto del sistema NutriRAG'' generada por Gemini.}
  \label{fig:c4}
\end{figure}

% ============================================================
\subsection{Trabajos relacionados}
% ============================================================

Las investigaciones recientes en sistemas de gestión nutricional asistidos
por inteligencia artificial evidencian una convergencia hacia arquitecturas
RAG con modelos de lenguaje locales y herramientas digitales
especializadas. La Tabla~\ref{tab:trabajos_relacionados} sintetiza los
trabajos de mayor relación metodológica y temática con el presente
proyecto; a continuación se presenta un análisis crítico de cada uno.

\begin{table}[ht]
\centering
\caption{Trabajos relacionados}
\label{tab:trabajos_relacionados}
\begin{tabular}{|c|p{9cm}|c|}
\hline
\textbf{ID} & \textbf{Título} & \textbf{Ref.} \\
\hline
TR01 & AI-Driven Personalized Nutrition: RAG-Based Digital Health Solution & \cite{gavai2025} \\
\hline
TR02 & Personalized Nutrition in the Era of Digital Health & \cite{arshad2025} \\
\hline
TR03 & Integrated Digital Health Care Platform for Diabetes Management & \cite{lee2023} \\
\hline
TR04 & Management and Monitoring of Patients in their Diet's Nutrition & \cite{quichimbo2021} \\
\hline
TR05 & Digital Applications for Diet Monitoring and Precision Nutrition & \cite{abeltino2025} \\
\hline
\end{tabular}
\end{table}

\textbf{TR01:} En el artículo~\cite{gavai2025}, los autores proponen un sistema de
nutrición personalizada basado en~RAG con modelos de lenguaje locales
orientado al manejo de la obesidad y la diabetes tipo~2. La solución
integra datos clínicos del paciente con una base de conocimiento
nutricional estructurada, reportando hasta un~80\% de adherencia y
capacidad de ajuste dinámico del plan. Su limitación principal es la
validación clínica aún incipiente, lo que justifica la inclusión de una
evaluación de satisfacción del profesional en el presente proyecto.

\medskip

\textbf{TR02:} En el artículo~\cite{arshad2025}, los autores presentan una revisión
de las tecnologías de salud digital para la nutrición personalizada en el
manejo de la diabetes y la obesidad, destacando el papel de los~LLM y los
sistemas~RAG en la generación de recomendaciones dietéticas adaptativas.
Los autores subrayan que integrar datos antropométricos y bioquímicos con
catálogos nutricionales estructurados mejora la pertinencia clínica de las
recomendaciones, a la vez que advierten sobre la necesidad de mantener la
supervisión del profesional como validador final del plan.

\medskip

\textbf{TR03:} Los autores del artículo~\cite{lee2023} evaluaron una plataforma
digital integrada para la gestión de la diabetes con módulos de
inteligencia artificial y gestión dietética durante 48~semanas. Los
resultados demuestran mejoras en los marcadores clínicos, pero ponen en
evidencia que los errores en la predicción de macronutrientes siguen siendo
un desafío sin supervisión profesional~\cite{mortazavi2023}. Esta
limitación refuerza la pertinencia de orientar el presente proyecto hacia
la satisfacción del nutricionista como criterio de aceptación del sistema.

\medskip

\textbf{TR04:} En el artículo~\cite{quichimbo2021}, los autores desarrollaron una
solución web para la gestión y monitoreo nutricional de pacientes en el
contexto ecuatoriano. Su trabajo constituye el antecedente nacional más
directo e identifica las mismas limitaciones que motivan el presente
proyecto: ausencia de herramientas inteligentes para generar
recomendaciones personalizadas y dependencia de registros manuales.
El prototipo propuesto avanza sobre estas limitaciones al incorporar un
pipeline~RAG con un catálogo de alimentos ecuatorianos.

\medskip

\textbf{TR05:} En el artículo~\cite{abeltino2025}, los autores realizaron una
revisión sistemática de aplicaciones digitales para el monitoreo, la
planificación y la nutrición de precisión. Identifican brechas críticas:
falta de interoperabilidad, escasa adaptación regional y altos costos de
implementación. Los autores concluyen que se requieren soluciones adaptadas
al contexto local con interfaces diseñadas para profesionales de la salud,
lo que fundamenta directamente el alcance del prototipo aquí propuesto.

% ============================================================
\subsection{Aporte del marco teórico al proyecto}
% ============================================================

El marco teórico establece los fundamentos que justifican cada decisión de
diseño del prototipo. La revisión de la gestión nutricional clínica en el
contexto ecuatoriano confirma las limitaciones del modelo manual
predominante y sustenta la necesidad de una herramienta digital que asista
al nutricionista~\cite{vargas2025,quichimbo2021}. El enfoque de Nutrición
Personalizada orienta la integración de datos antropométricos, bioquímicos
y del catálogo de alimentos ecuatorianos como entradas del sistema,
priorizando la pertinencia clínica local~\cite{siddiqui2025,aldarras2025}.

La arquitectura RAG, respaldada por evidencia empírica en sistemas de
nutrición similares~\cite{gavai2025,arshad2025}, resuelve la limitación
de las alucinaciones de los LLM al anclar las recomendaciones en datos
verificados. El uso de modelos de código abierto~(Qwen, Kimi~K, GLM o
Nemotron) ejecutados localmente responde a las limitaciones de costo,
interoperabilidad y privacidad identificadas en la
literatura~\cite{abeltino2025,ebadi2025}. Finalmente, la escala Likert
como instrumento de medición permite cuantificar el nivel de aceptación
del sistema por parte del nutricionista, alineando la evaluación con las
mejores prácticas en sistemas de soporte a la decisión
clínica~\cite{ghose2021,lee2023}.
